\chapter{Large Sample Theory}

Sometimes it is difficult to determine the exact distribution of a statistic for a finite value 
of $n$, and the assumptions made about the probability distribution may not be valid. 
In such cases, if the limiting distributions of the statistics exist, then the problems of testing 
such hypothesis and of setting confidence intervals may be easily solved for large samples.
However, these methods will be only approximate for a finite value of $n$.

On the bright side, these large-sample approximate results can make advantage since 
we do not have to make assumptions about the parent population and that in most cases 
the limiting distribution is normal. Thus normal theory can be applied to get 
approximate tests and confidence intervals. Another important 
limiting distribution used in connection with categorical data is the chi-squared distribution.
Such large-sample methods are beneficial in practical applications.

If the symptotic distribution of a statistic is normal, then for a large $n$ we can 
perform approximate large-sample tests on mean or variance of the statistic. In this chapter, we shall consider 
the case of categorical data and discuss certain transformations of statistics that are widely used in 
large samples.

\section{Transformation of statistics to stabilize variance}

\begin{lemma}
    \label{lem:las01}
    If $\{T_n\}_{n \geq 1}$ is a sequence of statistics such that 
    \begin{equation}
        \sqrt{n}(T_n - \theta) \xrightarrow[]{\mathcal{D}} Z \sim N\left(0, \sigma^2_T(\theta) \right),
    \end{equation}
    then the transformation of statistics $g(T_n)$ formed another sequence that
    \begin{equation}
        \sqrt{n}\left[ g(T_n) - g(\theta)\right] \xrightarrow[]{\mathcal{D}} \widetilde{Z} \sim N\left(0, \left[\sigma^2_T(\theta) g'(\theta) \right]^2 \right),
    \end{equation}
    provided $g(\theta)$ is differentiable and is nonzero.
\end{lemma}
\begin{proof}
    Consider the Taylor expansion of $g(T_n)$ around $g(\theta)$, that is,
    \[
        g(T_n) = g(\theta) + (T_n - \theta) [g'(\theta) + \epsilon_n]
    \]
    where $\epsilon \to 0$ as $T_n \to \theta$. Therefore, for any $\epsilon > 0$, $|\epsilon_n|<\epsilon$ whenerver 
    $|T_n - \theta| < \delta$. Hence, for any $\epsilon$ greater than zero, we have 
    \begin{align*}
        \mathbb{P} \left[ |\epsilon_n| < \epsilon\right] &= \mathbb{P} \left[ |T_n - \theta| < \delta \right]\\
        &= \mathbb{P} \left[ \frac{\sqrt{n} |T_n - \theta|}{\sigma_T(\theta)} < \frac{\sqrt{n} \delta}{\sigma_T(\theta)}\right]\\
        &= \mathbb{P} \left[ |\tau| < \frac{\sqrt{n} \delta}{\sigma_T(\theta)}\right] & \text{where } \tau \sim N(0,1)\\
        &= \left\{2 \Phi \left( \frac{\sqrt{n} \delta}{\sigma_T(\theta)}\right) - 1 \right\} \to 1 \quad \text{as } n \to \infty.
    \end{align*}
    By definition, $\epsilon$ converges to zero in probability. So this yields
    \[
        \sqrt{n}\left[ g(T_n) - g(\theta)\right] - \sqrt{n}(T_n - \theta)g'(\theta) = \fbox{$\sqrt{n}(T_n - \theta)\, \epsilon \xrightarrow[]{p} 0$}.
    \]
    Therefore, 
    \[
        \sqrt{n}\left[ g(T_n) - g(\theta)\right] \xrightarrow[]{\mathcal{D}} \sqrt{n}(T_n - \theta)g'(\theta) \sim N\left(0, \left[\sigma^2_T(\theta) g'(\theta) \right] \right).
    \]
    This completes the proof.
\end{proof}

\section{Stabilization of variance}

\index{Variance stabilization} is a statistical technique to 
make the spread (variance) of data consistent across its range by removing its dependence on the mean, which simplifies 
analysis and improves model performance by creating homoscedasticity (constant variance). It is crucial for data like 
counts (Poisson, Binomial) where variance naturally changes with the average, allowing 
standard methods (like ANOVA, regression) to work effectively and yield more reliable results, especially for small samples. 

In many real-world datasets (such as 
\textit{counts of rare events}, \textit{proportions}), the variability (variance) increases as the average value (mean) increases, 
violating assumptions of many statistical tests.

\begin{theorem}[Variance Stabilizing Transformation]
    If $\{T_n\}_{n \geq 1}$ is a sequence of statistics such that $\sqrt{n}(T_n - \theta) \sim N\left(0, \sigma^2_T(\theta) \right)$, 
    where $\sigma^2_T(\theta)$ is a function depends on $\theta$. Then 
    \begin{equation}
        g(\theta) \varpropto \int \frac{\mathrm{d}\theta}{\sigma_T(\theta)}
    \end{equation}
    is a \textbf{variance stabilizing transformation}.
\end{theorem}
\begin{proof}
    From \cref{lem:las01}, we have 
    \[
        \sqrt{n}\left[ g(T_n) - g(\theta)\right] \xrightarrow[]{\mathcal{D}} \widetilde{Z} \sim N\left(0, \left[\sigma^2_T(\theta) g'(\theta) \right]^2 \right)
    \]
    provided $g'(\theta)$ exists. We wish to find a transformation $g(T_n)$ whose variance is independent of $\theta$ (or, the variance is stable with respect to $\theta$). 
To stabilize the variance, we pick a function $g(T_n)$, such that 
\[
    [\sigma_T(\theta)\, g'(\theta)]^2 = c,
\]
where $c$ is a constant that is independent of $\theta$. The asymptotic variance of the transformed statistic $g(T_n)$ will be independent of $\theta$. 

Now, solving the differential equation $\sigma_T(\theta)\, g'(\theta) = c$.
\begin{align*}
    \sigma_T(\theta)\, g'(\theta) = c &\implies g'(\theta) = \frac{c}{\sigma_T(\theta)}\\
    &\implies \int g'(\theta) \> \mathrm{d}\theta = c \int \frac{\mathrm{d}\theta}{\sigma_T(\theta)}\\
    &\implies g(\theta) = c \int \frac{\mathrm{d}\theta}{\sigma_T(\theta)}.
\end{align*}
$g(\theta)$ is now proportional to the integral of the reciprocal of standard deviation.
\end{proof}

From the proof we know that $\sqrt{n}\left[ g(T_n) - g(\theta)\right]$ is asymptotically normal with mean $0$ and 
variance $c^2$, and that 
\begin{equation}
    \frac{\sqrt{n}\left[ g(T_n) - g(\theta)\right]}{c} \sim AN(0,1).
\end{equation}
When testing $H_0 : \theta = \theta_0$ versus $H_0 : \theta \neq \theta_0$ at significance
level $\alpha$, we reject $H_0$ if 
\begin{equation}
    \left| \frac{\sqrt{n}\left[ g(T_n) - g(\theta)\right]}{c} \right| > z_{\alpha/2}.
\end{equation}

\subsubsection{Arcsine Square Root Transformation}